\chapter{Introducción}
\label{cap:introduccion}

\section{Motivación}
\label{sec:motivacion}

Quien espera un autobús en una parada no necesita saber a qué hora \emph{debería} pasar, sino cuánto falta para que
pase \emph{de verdad}. La diferencia entre ambas cosas es el objeto de esta memoria. En una red interurbana como
Bizkaibus, con líneas que combinan calles congestionadas, carreteras y autopistas, esa diferencia puede ser de muchos
minutos y cambia con la hora del día, el día de la semana y el tramo concreto del recorrido.

\subsection{Esperar un Autobús Interurbano}
\label{sec:esperar}

En una línea urbana que pasa cada cinco minutos, el viajero no necesita saber a qué hora llega el autobús: va a la
parada y espera, en promedio, la mitad del intervalo. En una red interurbana eso no funciona, porque los intervalos son
largos. En el horario de Bizkaibus de un miércoles laborable (28 de octubre de 2026), entre las 07:00 y las 21:00, la
mitad de las combinaciones de línea y parada tienen, como mucho, un autobús por hora: el intervalo mediano es de
60~minutos, el 68\,\% de las combinaciones tiene un intervalo mediano de una hora o más, el 87\,\% de media hora o más,
y solo el 3\,\% tiene un autobús cada cuarto de hora o con más frecuencia.\footnote{Calculado sobre las 4\,932
combinaciones de línea y parada con servicio ese día, tomando para cada una el intervalo mediano entre salidas
consecutivas de la misma línea en la misma parada.}

Con esos intervalos, el viajero planifica: decide a qué hora salir de casa a partir de la hora a la que cree que
pasará el autobús. Si llega tarde, aunque sea por un minuto, pierde el autobús y espera el siguiente, a menudo una
hora. Si llega con mucho margen, espera en una parada que con frecuencia está en el arcén de una carretera, sin
marquesina ni información. Cuanto menos fiable es la hora de que dispone, más margen deja y más tiempo pierde.

\subsection{Lo que Cuesta No Saber Cuándo Llega}
\label{sec:coste-incertidumbre}

La literatura sobre transporte público ha medido el efecto de esta incertidumbre. Los viajeros que esperan sin
información perciben su espera como más larga de lo que es \citep{mishalani2006}. En Seattle, los usuarios de una
aplicación con llegadas en tiempo real no percibían su espera como más larga que la medida, a diferencia de quienes no
la usaban, y además esperaban menos \citep{watkins2011}. Una revisión de la literatura sobre los beneficios de la
información en tiempo real recoge, entre otros, menos tiempo de espera, menos tiempo total de viaje, más satisfacción,
más sensación de seguridad y, en algunos estudios, más uso del transporte público \citep{brakewood2019}.

Estos beneficios dependen de que la información sea \emph{precisa}. Una hora equivocada no solo no ayuda: hace perder
el autobús o esperar más que sin ella.

\subsection{El Horario Publicado Se Desvía en los Dos Sentidos}
\label{sec:desvio-horario}

Las horas de paso por las paradas intermedias que publica el horario de Bizkaibus son, como se explica en la
sección «\nameref{sec:problema}», una interpolación a velocidad constante entre cabeceras. Son las que muestra cualquier
aplicación que use el horario oficial en formato GTFS sin otra corrección. Para saber cuánto se desvían de la realidad,
se compararon con el paso real de los autobuses por las paradas, reconstruido a partir de su GPS
(sección «\nameref{sec:paso-real}»): 169\,308 pasos por parada observados entre el 6 y el 31 de octubre de 2026,
frente al horario vigente el 30 de octubre.\footnote{Se descartan 521 pasos (el 0,3\,\%) con más de 45~minutos de
diferencia, que suelen indicar un cambio de horario entre la fecha de la observación y la del horario usado. El cálculo
está en \texttt{referencia/desvio-horario.py}.} La Tabla~\ref{tab:desvio-horario} y la Figura~\ref{fig:desvio-horario}
muestran el resultado según el tiempo transcurrido desde la salida de cabecera.

\begin{table}[ht]
\centering
\caption{Desvío del Paso Real Respecto a la Hora Publicada, Según el Tiempo Desde la Salida}\label{tab:desvio-horario}
\begin{tabular}{@{}lrrrrr@{}}
\toprule
\shortstack[l]{Desde la\\cabecera} & Pasos & \shortstack{Desvío\\mediano} & \shortstack{Más de 2 min\\antes} & \shortstack{Más de 5 min\\después} & \shortstack{Dentro de\\$\pm2$ min}\\
\midrule
0--10 min   & 29\,115 & $+1{,}3$ min  & 5,0\,\%  & 11,8\,\% & 57,0\,\%\\
10--20 min  & 27\,408 & $-0{,}2$ min  & 33,3\,\% & 15,0\,\% & 34,9\,\%\\
20--30 min  & 28\,449 & $-2{,}2$ min  & 51,5\,\% & 14,4\,\% & 23,6\,\%\\
30--45 min  & 35\,560 & $-5{,}6$ min  & 67,6\,\% & 10,5\,\% & 14,6\,\%\\
45--60 min  & 25\,931 & $-8{,}6$ min  & 74,8\,\% & 7,4\,\%  & 11,6\,\%\\
60 min o más & 22\,845 & $-13{,}6$ min & 85,4\,\% & 3,7\,\%  & 8,0\,\%\\
\midrule
Todos       & 169\,308 & $-2{,}4$ min & 52,1\,\% & 10,7\,\% & 25,3\,\%\\
\bottomrule
\end{tabular}
\nota{El tiempo se cuenta en minutos programados desde la salida de cabecera. Un desvío negativo significa que el autobús pasa \emph{antes} de la hora publicada.}
\end{table}

\begin{figure}[ht]
\centering
\caption{Desvío del Paso Real Según el Tiempo Programado Desde la Salida}\label{fig:desvio-horario}
\begin{tikzpicture}
\begin{axis}[
  width=0.92\textwidth, height=6cm,
  ybar, bar width=9pt,
  symbolic x coords={0--10,10--20,20--30,30--45,45--60,60+},
  xtick=data,
  xlabel={Minutos programados desde la salida de cabecera},
  ylabel={Pasos por parada (\%)},
  ymin=0, ymax=100,
  axis lines*=left, ymajorgrids, grid style={gris!20},
  tick label style={font=\footnotesize}, label style={font=\small},
  legend style={font=\footnotesize,draw=none,fill=none,at={(0.5,-0.28)},anchor=north,legend columns=2,/tikz/every even column/.append style={column sep=1em}},
  nodes near coords, every node near coord/.append style={font=\scriptsize,color=black},
  nodes near coords style={/pgf/number format/fixed,/pgf/number format/precision=0}]
\addplot[fill=gps,draw=gps] coordinates {(0--10,5.0) (10--20,33.3) (20--30,51.5) (30--45,67.6) (45--60,74.8) (60+,85.4)};
\addlegendentry{Pasa más de 2 min antes de la hora publicada}
\addplot[fill=ruta,draw=ruta] coordinates {(0--10,11.8) (10--20,15.0) (20--30,14.4) (30--45,10.5) (45--60,7.4) (60+,3.7)};
\addlegendentry{Pasa más de 5 min después}
\end{axis}
\end{tikzpicture}
\nota{Cerca de la cabecera el autobús va algo retrasado; a medida que avanza, se adelanta cada vez más a la hora publicada, porque el horario reparte el tiempo a velocidad constante y el autobús recorre los tramos de carretera más deprisa.}
\end{figure}

El resultado es el que predice la interpolación. Cerca de la cabecera el autobús va algo retrasado (desvío mediano de
$+1{,}3$~min en los diez primeros minutos). A partir de ahí se adelanta cada vez más al horario: a media hora de la
salida, la mitad de los pasos ocurre más de dos minutos antes de la hora publicada, y pasada la hora, el 85\,\%, con un
adelanto mediano de casi catorce minutos. En fin de semana el adelanto es mayor: se adelantan más de dos minutos el
60\,\% de los pasos, frente al 50\,\% en laborables. De las 69 líneas con al menos 500 pasos observados, en 39 la
mayoría de los pasos se adelanta más de dos minutos. Solo una cuarta parte de los pasos cae dentro de $\pm2$~minutos de la hora publicada.

Para el viajero, eso significa dos riesgos de signo contrario:

\begin{itemize}
  \item \textbf{Perder el autobús.} Quien se guiara por la hora publicada y llegara a la parada dos minutos antes
        encontraría que el autobús ya había pasado en el 52\,\% de los casos. Con un intervalo típico de una hora, cada
        autobús perdido es una hora de espera.
  \item \textbf{Esperar de más.} Cuando el autobús todavía no ha pasado, esa misma persona espera una mediana de
        3,8~minutos, pero una de cada diez veces, más de diez minutos, y el 10,7\,\% de los pasos llega más de cinco
        minutos tarde.
\end{itemize}

\noindent Estas cifras no son un juicio sobre la puntualidad del servicio: las horas intermedias del GTFS no son un
compromiso del operador (sección «\nameref{sec:problema}»), sino una aproximación. Son, sin embargo, las horas que ve quien
consulta el horario en una aplicación que no las corrige, y por eso el problema merece una solución.

\subsection{Este Trabajo}

Bizkaibus+ es una aplicación web independiente que muestra, para cada parada, las próximas salidas con su hora
estimada. Este documento describe el algoritmo que calcula esas horas: de dónde salen los datos, qué se hace con
ellos, por qué se hace así y cómo se ha comprobado que funciona. El objetivo es que el método pueda estudiarse,
criticarse y replicarse en cualquier otra ciudad y en cualquier lenguaje de programación.

\section{El Problema}
\label{sec:problema}

Los datos disponibles condicionan todo el diseño. Bizkaibus publica dos fuentes abiertas, que son las únicas que usa
este trabajo:

\begin{enumerate}
  \item Un \textbf{horario programado} en formato GTFS \citep{gtfs}, publicado en el Punto de Acceso Nacional de
        transporte multimodal \citep{nap}: líneas, viajes, paradas y la hora de paso de cada viaje por cada parada,
        junto con el trazado geográfico (\emph{shape}) de cada recorrido.
  \item Un \textbf{feed en tiempo real} en formato SIRI Vehicle Monitoring \citep{siri}, actualizado aproximadamente
        una vez por minuto, que para cada autobús en servicio indica el viaje que realiza, su posición GPS, la hora a la que se tomó
        esa posición y el número de orden de su próxima parada.
\end{enumerate}

Dos características de estas fuentes descartan el enfoque habitual de sumar un retraso al horario:

\begin{itemize}
  \item \textbf{El feed no publica retrasos.} El elemento \texttt{Delay} de SIRI no aparece, y el feed alternativo de
        actualizaciones de viaje declara siempre un retraso nulo (\texttt{PT0S}) con una marca de tiempo inválida.
  \item \textbf{El horario intermedio es una interpolación.} Las paradas intermedias del GTFS llevan
        \texttt{timepoint=0}: sus horas se obtienen repartiendo el tiempo entre cabeceras a velocidad constante. Un
        autobús que recorre 18~km de autopista en 15~minutos puede tener, para ese tramo, entre 34 y 49~minutos
        programados. En mitad del recorrido, el horario absoluto puede ir quince o veinte minutos por detrás del
        autobús sin que este se haya adelantado en ningún sentido operativo.
\end{itemize}

Así, ni existe un retraso que sumar ni el horario absoluto es una referencia fiable. Lo único vivo y fiable es la
posición.

\section{Objetivos}

\begin{enumerate}
  \item Estimar, para cada viaje en servicio y cada parada que le quede por recorrer, la hora de llegada con el menor
        error posible, a partir de la posición del vehículo.
  \item Seguir dando una estimación útil cuando el GPS falla o no existe, sin inventar datos.
  \item Aprender de forma automática el comportamiento real de la red, sin reglas escritas a mano sobre la velocidad
        de los autobuses.
  \item Medir la precisión contra el paso real de los autobuses por las paradas, y adoptar un cambio en el modelo solo
        si esa medida mejora.
  \item Documentarlo todo con la precisión suficiente para que otra persona pueda reconstruirlo sin leer el código.
\end{enumerate}

\section{Contribuciones}

\begin{itemize}
  \item Un método de predicción por \textbf{diferencias de horario} que neutraliza el error de forma de un horario
        interpolado (capítulo «\nameref{cap:modelo}», proposición~\ref{prop:diferencias}).
  \item Un procedimiento de \textbf{localización} que sitúa el GPS sobre el trazado oficial y es robusto ante los
        recorridos que vuelven sobre sí mismos (capítulo «\nameref{cap:localizacion}»).
  \item Un \textbf{modelo de tiempos por tramo} con estimadores robustos, contracción hacia la mediana del tramo y una
        cascada de respaldo para los tramos sin datos (capítulo «\nameref{cap:modelo}»).
  \item Una arquitectura de \textbf{tres capas con abstención}, que elige en cada predicción la mejor fuente fiable y
        nunca afirma un retraso que no pueda sostener, con una \textbf{corrección por viaje} para cuando no hay ningún
        dato del vehículo (capítulo «\nameref{cap:prediccion}»).
  \item Una \textbf{revisión de los datos de entrada} en la construcción del horario: los trazados dibujados al revés se
        invierten y, en las redes que no publican los metros de cada parada, estos se miden sobre el trazado
        (capítulo «\nameref{cap:datos}»).
  \item Un ciclo de \textbf{aprendizaje continuo} con validación cruzada y reglas de escritura conservadoras
        (capítulo «\nameref{cap:aprendizaje}»), un protocolo de \textbf{evaluación} por días que reproduce lo que ve el
        usuario, y una \textbf{comparación} con filtros de Kalman y redes neuronales hecha con ese mismo protocolo
        (capítulo «\nameref{cap:evaluacion}»).
  \item \textbf{Vectores de prueba} generados con el código de producción y una \textbf{implementación de referencia}
        independiente que los reproduce exactamente (Apéndices~\ref{ap:vectores} y~\ref{ap:referencia}).
\end{itemize}

\section{Visión General del Sistema}

La Figura~\ref{fig:tuberia} resume el recorrido de un dato desde el feed hasta la pantalla. Cada paso se describe en
el capítulo correspondiente.

\begin{figure}[ht]
\centering
\caption{Recorrido de un Dato Desde el Feed Hasta la Pantalla}\label{fig:tuberia}
\begin{tikzpicture}[
  caja/.style={draw=ruta,rounded corners=3pt,thick,align=center,minimum height=1.25cm,text width=2.25cm,font=\footnotesize},
  flecha/.style={-{Stealth[length=2.2mm]},thick,color=ruta}]
  \node[caja] (feed) {\textbf{Feed SIRI-VM}\\posición, próxima parada, hora};
  \node[caja,right=0.45cm of feed] (viaje) {\textbf{Emparejar}\\con el viaje del GTFS};
  \node[caja,right=0.45cm of viaje] (ruta) {\textbf{Proyectar}\\el GPS en el trazado ($\Tpos$)};
  \node[caja,right=0.45cm of ruta] (tramos) {\textbf{Sumar tramos}\\tiempo aprendido hasta la parada};
  \node[caja,below=0.9cm of tramos] (eta) {\textbf{Llegada}\\estado y retraso};
  \node[caja,below=0.9cm of feed] (captura) {\textbf{Captura}\\continua, por horas};
  \node[caja,right=0.45cm of captura] (muestras) {\textbf{Muestras}\\tramos y ritmo};
  \node[caja,right=0.45cm of muestras] (calibrar) {\textbf{Recalibrar}\\cada día, si mejora};
  \draw[flecha] (feed) -- (viaje);
  \draw[flecha] (viaje) -- (ruta);
  \draw[flecha] (ruta) -- (tramos);
  \draw[flecha] (tramos) -- (eta);
  \draw[flecha] (captura) -- (muestras);
  \draw[flecha] (muestras) -- (calibrar);
  \draw[flecha,dashed] (calibrar.north east) -- (tramos.south west);
\end{tikzpicture}
\nota{Arriba, la predicción en cada petición (capítulos «\nameref{cap:localizacion}» a «\nameref{cap:prediccion}»). Abajo, el aprendizaje fuera de línea que alimenta los tiempos por tramo (capítulo «\nameref{cap:aprendizaje}»).}
\end{figure}

\section{Convenciones}
\label{sec:convenciones}

\begin{itemize}
  \item \textbf{Tiempo.} Todas las horas son números enteros de segundos desde la medianoche de la zona
        \texttt{Europe/Madrid} del día de servicio. Siguiendo a GTFS, pueden superar $86\,400$: las 25:10 son
        $90\,600$~s. El feed y el reloj dan horas del día en curso; para un viaje del día anterior se les suman
        $86\,400$~s (sección «\nameref{sec:mas-de-24h}»). Las duraciones se expresan en segundos salvo que se diga otra
        cosa.
  \item \textbf{Distancia.} En metros. Las coordenadas geográficas, en grados decimales (WGS84).
  \item \textbf{Redondeo.} $\redondeo(x)$ redondea al entero más próximo con las mitades \emph{alejándose del cero}
        ($\redondeo(2{,}5)=3$, $\redondeo(-2{,}5)=-3$). Es el comportamiento de PHP y difiere del «redondeo del
        banquero» de otros lenguajes; replicarlo es necesario para obtener los mismos resultados.
  \item \textbf{Mediana.} Salvo que se indique otra cosa, $\mediana$ es la \emph{mediana baja}: con los valores
        ordenados $x_{(0)}\le\dots\le x_{(n-1)}$, $\mediana(x)=x_{(\lfloor (n-1)/2\rfloor)}$. En general, el percentil
        $p$ es $P_p(x)=x_{(\lfloor p\,(n-1)\rfloor)}$, sin interpolar.
  \item \textbf{Acotación.} $\clamp(x,a,b)=\max(a,\min(b,x))$.
  \item \textbf{Nombres en el código.} Cuando conviene, se indica entre paréntesis y en letra \texttt{monoespaciada} el
        nombre de la constante o función del código, para facilitar la correspondencia.
\end{itemize}

\section{Estructura del Documento}

El documento sigue la estructura habitual de un trabajo de investigación (introducción, método, resultados y
discusión) y el formato de las normas APA en su séptima edición. Este capítulo y el capítulo «\nameref{cap:marco}»
sitúan el trabajo respecto a la literatura de predicción de llegadas y a las herramientas estadísticas que usa. Los
cinco siguientes describen el método completo: las fuentes de datos (capítulo «\nameref{cap:datos}»), la localización del
vehículo (capítulo «\nameref{cap:localizacion}»), el modelo del tiempo de recorrido (capítulo «\nameref{cap:modelo}»), la
predicción en tres capas (capítulo «\nameref{cap:prediccion}») y el aprendizaje de los parámetros
(capítulo «\nameref{cap:aprendizaje}»). El capítulo «\nameref{cap:evaluacion}» presenta la precisión medida en una muestra
con las tres capas, a lo largo de 17 días y frente a otros métodos; el capítulo «\nameref{cap:discusion}» la interpreta y
discute las limitaciones del método y su aplicación a otras redes, y el capítulo «\nameref{cap:conclusiones}» cierra con las
conclusiones y las líneas futuras. Tras las referencias, los apéndices reúnen las constantes, el pseudocódigo
completo, los formatos de datos, los vectores de prueba, la implementación de referencia y un glosario.
